\chapter{Background and Related Work}
\label{chap:background}

\section{Knowledge graphs and scholarly records}

A knowledge graph represents entities and typed relationships in a form that can
be queried and combined. In this project, the graph is a property graph: domain
nodes have stable identifiers and typed labels, and an edge represents one
predicate between two nodes. The design is intentionally small enough for a
course analysis but detailed enough to separate papers, people, organizations,
venues, and inferred topics. Surveys of knowledge-graph systems emphasize that
representation, construction, quality assessment, and application are connected
problems \citep{hogan2022knowledgegraphs}. Here, each relationship also carries
an evidence class and links to source or support records so that the graph can be
traced back to its construction inputs.

Scholarly metadata is heterogeneous even when records describe the same paper.
Proceedings pages and index records can differ in title formatting, author
coverage, identifiers, abstracts, or affiliation detail. Database-level
affiliation discrepancies have been documented across bibliographic sources
\citep{purnell2022affiliation}. The project's response is to assign different
roles to sources: official proceedings define inclusion and byline order, while
OpenAlex is used for work metadata, author IDs, institutions, and reference lists.
The second source enriches the record but does not replace the official corpus
boundary.

\section{Scientific collaboration networks}

Coauthorship networks model researchers as vertices and shared publications as
evidence of collaboration. Newman showed how scientific collaboration networks
can be studied through degree, paths, components, and clustering
\citep{newman2001collaboration}. The meaning of an edge nevertheless depends on
the counting rule. A pair may have one shared paper or many; a weighted edge can
retain the number or fractional contribution of those papers. This report
therefore distinguishes unique coauthor neighbors from fractional strength and
reports the accepted graph's support rules.

The project also constructs organization collaboration from paper-specific
affiliations. Two institutions are connected only when at least two distinct
authors on a core paper have accepted affiliations to the two organizations.
A single researcher listing two organizations does not establish an
organization-to-organization collaboration. This rule avoids interpreting a
dual-affiliation record as evidence of a joint paper between institutions.

\section{Entity resolution and identity uncertainty}

Entity resolution links records that refer to the same real-world object. In
bibliographic settings, a person can publish under name variants, while two
different people can share the same name. Record-linkage methods therefore use
multiple identifying attributes and explicit decisions instead of relying on a
single normalized string \citep{christen2012datamatching}. A provider's persistent
author identifier improves record linkage but is still a provider assertion;
conflicting identifiers or incompatible author names require abstention or
review.

This project uses OpenAlex author IDs at compatible byline positions and stores
ORCID when present. It does not merge researchers by name similarity alone.
The corrected v4 graph also applies a global conflict rule to observed ORCID
evidence across all matched works. This makes the identity policy conservative
and reversible, but does not independently prove that every surviving provider
ID belongs to the printed person. Downstream coauthor and bridge claims inherit
that uncertainty.

\section{Text classification and topic assignment}

Sentence embeddings represent text in a vector space where semantic similarity
can be compared. Sentence-BERT introduced a siamese architecture for producing
useful sentence-level embeddings \citep{reimers2019sentencebert}. This project
uses a frozen MiniLM sentence-transformer to compare paper title/abstract text
with definitions in a small curated taxonomy. The classifier is a fixed
embedding-and-threshold procedure, not an open-ended large-language-model
judgment. A paper receives between one and three leaf topics or no topic label
when the confidence rule abstains.

The distinction between coverage and correctness is essential. Coverage records
the proportion of papers assigned at least one label. Precision, recall, and F1
compare assigned labels with reference labels and depend on the provenance and
quality of those references. In this project, the available reference labels are
AI-assessed, drawn from exposed frozen test papers, and are not independent human
gold. Topic portfolios and trends are consequently descriptive and provisional.

\section{Network measures and community detection}

Degree counts distinct adjacent nodes; weighted degree, or strength, sums edge
weights. Betweenness centrality measures how often a vertex lies on shortest
paths between other vertices, under a specified path length and graph projection
\citep{freeman1977betweenness}. For collaboration analysis, this project treats
the inverse of positive fractional strength as edge distance. High betweenness
therefore identifies a potential structural bridge under that particular
distance definition, not an individual's research quality or causal influence.

Community detection groups nodes by comparatively dense internal connections.
The release uses a configured Louvain method run with fixed seed and resolution;
the method optimizes modularity through multilevel community aggregation
\citep{blondel2008communities}. A community assignment is conditional on the
projection, edge weights, seed, and resolution. It should not be treated as an
official disciplinary boundary or a stable identity label across graph versions.

\section{Evidence provenance and graph quality}

Graph quality has structural and semantic dimensions. Structural checks can
confirm unique identifiers, valid endpoints, allowed labels and predicates,
evidence references, support counts, and database reconciliation. They cannot
establish that a source record is correct or that a provider identity matches a
real person. The project therefore records evidence URLs, retrieval timestamps,
raw-response hashes, locators, and method/graph versions for accepted edges.
Derived edges point to supporting base edges and papers, while ambiguous claims
remain outside the accepted graph.

The edge review reports precision for sampled accepted edges; it does not sample
missing relationships and therefore does not estimate recall. For the principal
AI-assessed audit, a Wilson interval summarizes binomial sampling uncertainty
conditional on the supplied judgments \citep{wilson1927inference}. It does not
include evaluator mistakes, shared-model bias, or dependence on construction
sources. This narrow interpretation is retained throughout the results.

\section{How this project differs from downloading a graph}

The project defines its own corpus from official proceedings rosters, performs
DOI-based record matching, applies explicit byline and affiliation rules, creates
topic links from a frozen classifier, derives four social and topic relationships
from accepted support, and validates the graph's internal evidence structure.
Existing OpenAlex records supply metadata but are not used as a prebuilt graph
whose edges are accepted without inspection. The resulting graph is thus a
versioned analytical artifact with a traceable evidence contract, not a direct
copy of a third-party knowledge graph.
